\begin{exercise}[Gauge-invariant fields in the unitary-gauge limit]{I-CH08-020}
Let a scalar vacuum $v$ break $G$ to $H$, and decompose fields by
\[
 \phi=\Omega(v+\Delta),\qquad
 B_\mu=\Omega^\dagger(\partial_\mu-gA_\mu)\Omega,
 \qquad \Psi=\Omega_F^\dagger\psi,
\]
where $\Delta$ is orthogonal to the gauge orbit of $v$.  Expand
$\Omega=\exp\pi$ through first order and derive the transformation laws of
$B_\mu$, $\Delta$, and $\Psi$, including the compensating action of the
unbroken subgroup $H$.
Identify the dressed broken-vector combination, the physical Higgs
directions, and the dressed fermion.  In an $R_\kappa$ gauge, use the
$\kappa$-dependent Goldstone, ghost, and unphysical-vector masses to show that
the $\kappa\to\infty$ limit sets the gauge-orbit fluctuation to zero in
low-energy pole residues.  Demonstrate that the surviving massive fields
then agree with these combinations to leading order, with residual
$H$ covariance when $H$ is nontrivial.
Explain why vectors in the unbroken subgroup $H$ still require a transverse
projection and why this limit cannot be interchanged with the ultraviolet
limit of general Green functions.
\end{exercise}
